% Requires in preamble: \usepackage{booktabs, amsmath, cleveref}
\subsection{Energy Performance}
The Energy Performance Certificate (EPC) is an official document that assesses a
building's energy consumption. It is mandatory for every sale or rental in
Belgium and assigns a label ranging from A (very efficient) to G (very
inefficient). The score indicates the theoretical amount of primary energy used
by the building for space heating and domestic hot water over one year, expressed
in kWh/m$^2$, and depends on the insulation of the walls, roof, floors, windows
and doors, on the heating installations, on the presence of solar panels, and on
characteristics such as the type of building (detached, semi-detached, terraced
or apartment) \cite{nbb_epc_houseprices_2025}. While the energy score is
calculated using very similar methodologies in all three Belgian regions, the
grouping of scores into energy labels differs markedly between regions and yields
different scales, as shown in \Cref{tab:epc_regional_scales}
\cite{nbb_epc_houseprices_2025}. The EPC therefore summarises, in a single
indicator, the envelope quality and technical systems that the archetypes of
\Cref{sec:size_compo} represent through their U-values, airtightness and
construction period.

\begin{table}[ht]
  \centering
  \caption{Grouping of EPC scores (kWh/m$^2$) into energy labels by region
           \cite{nbb_epc_houseprices_2025}. The Flemish scale does not define a
           class~G.}
  \label{tab:epc_regional_scales}
  \begin{tabular}{lccc}
    \toprule
    Label & Flemish Region & Walloon Region & Brussels-Capital Region \\
    \midrule
    A & $\leq 100$   & $\leq 85$   & $\leq 45$   \\
    B & 101--200     & 86--170     & 46--95      \\
    C & 201--300     & 171--255    & 96--150     \\
    D & 301--400     & 256--340    & 151--210    \\
    E & 401--500     & 341--425    & 211--275    \\
    F & $> 500$      & 426--510    & 276--345    \\
    G & ---          & $> 510$     & $> 345$     \\
    \bottomrule
  \end{tabular}
\end{table}

\subsubsection*{Regional EPC distribution}
The Belgian residential stock is dominated by average to low-performing energy
labels, but the distribution varies by region because of differences in building
age, dwelling type and urban density. A useful indicator is the average energy
score of dwellings sold on the secondary market, which the National Bank of
Belgium reports for the first quarter of 2024 as 363, 390 and 364~kWh/m$^2$ for
houses, and 198, 230 and 245~kWh/m$^2$ for apartments, in the Flemish, Walloon
and Brussels-Capital Regions respectively \cite{nbb_epc_houseprices_2025}. These
figures describe \emph{sold} dwellings, whose energy performance is somewhat
worse than that of the full stock because new builds are excluded and sold homes
tend to be older; they nonetheless remain far from the 2050 targets of
100~kWh/m$^2$ (Flanders and Brussels) and 85~kWh/m$^2$ (Wallonia)
\cite{nbb_epc_houseprices_2025}.

\paragraph{Brussels-Capital Region.} Brussels combines a dense stock of older,
historic apartment buildings with the strictest label scale in the country (a
class~G already above 345~kWh/m$^2$; see \Cref{tab:epc_regional_scales}), so its
stock is concentrated in the lower classes. In a sample of 2{,}631 residential
certificates issued by a single Brussels certifier between 2011 and 2026, labels
F and G together account for about 54\% of dwellings, with class~G alone at
roughly 36\%; because this is a single-certifier sample rather than a
stock-representative one, these shares should be read as indicative
\cite{certibru_epc_stats}.

\paragraph{Flemish Region.} Flanders shows a comparatively more balanced
distribution, with higher shares of A, B and C labels driven by stricter
new-construction standards and a renovation obligation for the least efficient
homes. A substantial share of older single-family houses nonetheless remains in
the E and F classes: about 35\% of houses sold in the Flemish Region between 2023
and the first quarter of 2024 carried a class~E or~F label
\cite{nbb_epc_houseprices_2025}.

\paragraph{Walloon Region.} Wallonia has a large share of older, rural and
industrial-era single-family homes, which gives it the highest average energy
score of the three regions: sold Walloon houses averaged 390~kWh/m$^2$ in the
first quarter of 2024, corresponding to a class~E label on the Walloon scale
\cite{nbb_epc_houseprices_2025}.
